\chapter*{Introduction}
\addcontentsline{toc}{chapter}{Introduction}

\lettrine{A}{rrays are everywhere}. Every image you view, every text you read, every game you play, every database query you execute, every model that recognizes your voice---arrays underlie them all. Yet despite their ubiquity, arrays remain poorly understood by most programmers. They are treated as primitive constructs, learned hastily and used mechanically, their true nature obscured by layers of abstraction stacked between a line of code and the silicon that executes it.

\textit{Arliz} exists to remedy that situation. The Preface explained why I went looking for that remedy, what the name means, and why the work is divided into three volumes; this Introduction is about \emph{this} volume specifically---what it covers, what it assumes of you, and how it connects to the rest of the work.

\section*{This Volume in the Larger Work}

% NOTE: Each volume's introduction should adjust this section to describe

\textit{Arliz} follows one continuous path: from how information is encoded, to the machines that compute on that information, to the array itself in all its forms. That path is divided into three volumes:

\begin{itemize}
	\item \textbf{Volume I --- Zero to Bit} (this volume)\\
	We begin with the fundamental question: how is information encoded in digital systems at all? We start at the level of voltage and binary switching, then move through binary representation and Boolean algebra, number systems, integer and floating-point formats, character encoding, bitwise operations, alignment, and serialization. This volume establishes the representational substrate that determines how every array element is actually stored and manipulated at the machine level.

	\item \textbf{Volume II --- Architecture \& Circuitry}\\
	Picks up where this volume ends. Arrays do not exist in abstract space---they live in physical hardware with specific characteristics and constraints. Volume II examines semiconductor physics, logic gates, arithmetic circuits, memory cell design, the full memory hierarchy and cache behavior, processor pipelines, instruction set architectures, vector and SIMD units, GPU architecture, NUMA systems, and the interconnects that tie it all together.

	\item \textbf{Volume III --- Array Odyssey}\\
	The destination of the first two volumes. Arrays in full: their mathematics, memory layout, every major variant, the data structures and algorithms built on top of them, and the parallel and large-scale systems---from CPU threads to GPU kernels to distributed tensors---that process them today.
\end{itemize}

Each volume is written to be readable on its own, but the dependencies run in one direction: Volume II assumes the representational vocabulary built here in Volume I, and Volume III assumes both. If you find yourself in Volume III wondering why an operation costs what it does at the hardware level, or in Volume II wondering why a value is encoded the way it is, the earlier volume is where that question is answered.

\section*{What This Volume Covers}

This volume answers a single question, asked as literally as possible: how is information encoded in digital systems, starting from nothing but a voltage difference? We begin before computing altogether---with what data \emph{is}, and what it means for one thing to represent another---and then follow that question through binary switching, Boolean algebra, positional number systems, signed and unsigned integers, fixed- and floating-point arithmetic, character encoding and Unicode, endianness, bitwise manipulation, memory alignment, media representations (color, image, audio, video), pointers, error detection and correction, and serialization and compression. By the end of this volume, you will have every representational tool needed to understand, in Volume II, how hardware actually stores and moves these encodings, and, in Volume III, how arrays organize them into the structures that power real software.

\section*{Prerequisites}

This book assumes substantial preparation:

\textbf{Mathematical Maturity:} You should be comfortable with discrete mathematics, linear algebra, and basic mathematical analysis. Specifically, you should have completed (or be comfortable with) the material in \textit{Mathesis: The Mathematical Foundations of Computing}.

\textbf{Algorithmic Analysis:} You should understand asymptotic notation, recurrence relations, and basic complexity analysis. The material in \textit{The Art of Algorithmic Analysis} provides the necessary foundations.

\textbf{Programming Experience:} You should be proficient in at least one programming language and comfortable reading code in multiple languages. We use pseudocode, C, and occasionally other languages for examples.\\
Without these prerequisites, you will struggle with substantial portions of this volume. With them, you are prepared for a deep, rewarding start to the journey from voltage to vector.

\section*{The Living Nature of This Work}

Check the GitHub repository regularly for updates. The version you read today will be superseded by improved versions tomorrow---and because the work is now split across volumes, each one can be revised on its own schedule without forcing a reprint of the others. This is not a weakness but a strength: the work remains current, accurate, and relevant, and your corrections and questions are part of how it gets there.

\section*{A Note on Difficulty}

This volume is challenging by design, even though it is, in a sense, the most "elementary" of the three---it deals with representation rather than algorithms or hardware. But representation is where every later mistake is born: a misunderstood encoding produces bugs and performance problems that no amount of algorithmic cleverness can fix later. Some sections will require sustained effort to master.

This difficulty is intentional. Genuine understanding is never cheap---it demands intellectual investment, persistent effort, and willingness to struggle with complex concepts. If you find sections difficult, that is normal. Persist. Work through examples. Understanding will come through sustained engagement, and it will pay off directly when you reach Volumes II and III.

\section*{Begin}

This volume is the first stage of a longer journey---the one that starts at a voltage and, three volumes later, ends at a tensor. Everything here is foundation: necessary, occasionally dense, and built to be returned to.

Welcome to \textbf{Arliz}. Let us begin where everything begins---at the bit.

\vfill

\begin{quote}
\textit{``Premature optimization is the root of all evil.''}

\hfill--- \textsc{Donald E. Knuth}
\end{quote}

\clearpage